\documentclass[10pt,aspectratio=169]{beamer}

\usetheme{Madrid}
\usecolortheme{seahorse}
\setbeamertemplate{navigation symbols}{}
\setbeamertemplate{footline}{%
  \hfill\scriptsize\color{DMLMuted}\insertframenumber/\inserttotalframenumber\hspace{0.35cm}\vspace{0.18cm}
}

\usepackage[utf8]{inputenc}
\usepackage[T1]{fontenc}
\usepackage[spanish]{babel}
\usepackage{booktabs}
\usepackage{array}
\usepackage{graphicx}
\usepackage{tikz}
\usepackage{hyperref}

\definecolor{DMLInk}{RGB}{15,23,42}
\definecolor{DMLMuted}{RGB}{71,85,105}
\definecolor{DMLSoft}{RGB}{248,250,252}
\definecolor{DMLLine}{RGB}{203,213,225}
\definecolor{DMLGreen}{RGB}{22,101,52}
\definecolor{DMLBlue}{RGB}{29,78,216}
\definecolor{DMLBeamerBlue}{RGB}{35,74,145}
\definecolor{DMLGold}{RGB}{146,64,14}
\definecolor{DMLRed}{RGB}{153,27,27}

\setbeamercolor{normal text}{fg=DMLInk,bg=white}
\setbeamercolor{frametitle}{fg=white,bg=DMLBeamerBlue}
\setbeamercolor{title}{fg=DMLBeamerBlue}
\setbeamercolor{subtitle}{fg=DMLInk}
\setbeamercolor{structure}{fg=DMLBeamerBlue}
\setbeamercolor{palette primary}{fg=white,bg=DMLBeamerBlue}
\setbeamercolor{palette secondary}{fg=white,bg=DMLBlue}
\setbeamercolor{palette tertiary}{fg=white,bg=DMLBeamerBlue}
\setbeamercolor{block title}{fg=white,bg=DMLBeamerBlue}
\setbeamercolor{block body}{fg=DMLInk,bg=DMLSoft}
\setbeamerfont{frametitle}{series=\bfseries,size=\Large}
\setbeamerfont{title}{series=\bfseries,size=\Huge}
\setbeamerfont{subtitle}{size=\large}

\newcommand{\DMLEyebrow}[1]{%
  {\scriptsize\bfseries\color{DMLGreen}#1}\vspace{0.08cm}
}

\newcommand{\DMLBox}[1]{%
  \begin{tikzpicture}
    \node[
      draw=DMLLine,
      fill=DMLSoft,
      rounded corners=2pt,
      inner sep=8pt,
      text width=0.84\textwidth,
      align=left
    ] {\large #1};
  \end{tikzpicture}
}

\newcommand{\DMLSmallBox}[1]{%
  \begin{tikzpicture}
    \node[
      draw=DMLLine,
      fill=DMLSoft,
      rounded corners=2pt,
      inner sep=7pt,
      text width=0.86\linewidth,
      align=left
    ] {\small #1};
  \end{tikzpicture}
}

\newcommand{\DMLCard}[3]{%
  \begin{tikzpicture}
    \node[
      draw=#1!50,
      fill=#1!5,
      rounded corners=2pt,
      inner sep=7pt,
      text width=0.86\linewidth,
      minimum height=1.95cm,
      align=left
    ] {\textbf{\color{#1}#2}\par\vspace{0.08cm}{\small #3}};
  \end{tikzpicture}
}

\newcommand{\DMLMetric}[3]{%
  \begin{tikzpicture}
    \node[
      draw=#1!45,
      fill=#1!6,
      rounded corners=2pt,
      inner sep=8pt,
      text width=0.86\linewidth,
      minimum height=1.42cm,
      align=left
    ] {{\Large\bfseries\color{#1}#2}\par{\scriptsize #3}};
  \end{tikzpicture}
}

\newcommand{\DMLImage}[2]{%
  \IfFileExists{images/#1}{%
    \includegraphics[width=\linewidth,height=0.48\textheight,keepaspectratio]{images/#1}
  }{%
    \begin{tikzpicture}
      \node[
        draw=DMLLine,
        fill=DMLSoft,
        rounded corners=2pt,
        minimum width=\linewidth,
        minimum height=0.44\textheight,
        align=center,
        inner sep=10pt,
        text width=0.76\linewidth
      ] {\color{DMLMuted}\textbf{Imagen sugerida}\par\vspace{0.2cm}#2\par\vspace{0.25cm}{\scriptsize Generar y guardar como \texttt{\detokenize{#1}}}};
    \end{tikzpicture}
  }
}

\title{d-ml}
\subtitle{Digital Machine Listening aplicado al audio clínico}
\author{Mario Fernando Jojoa Acosta}
\date{Conversación estratégica con el Dr. Cantón · Mayo 2026}

\begin{document}

% 1
\begin{frame}[plain]
  \vspace{0.05cm}
  \begin{center}
    \DMLImage{01_sound_to_evidence.png}{Audio clínico transformándose en evidencia y datos.}
  \end{center}
  \vspace{-0.10cm}
  {\Huge\bfseries d-ml\par}
  \vspace{0.12cm}
  {\Large Digital Machine Listening: convierte sonidos clínicos —voz, pulmón y corazón— en datos accionables.\par}
  \vspace{0.22cm}
  {\small Mario Fernando Jojoa Acosta · Senior Research Engineer, Barcelona Supercomputing Center · \url{https://www.d-ml.eu}}
\end{frame}

% 2
\begin{frame}{Agenda de conversación}
  \DMLEyebrow{RUTA PROPUESTA · 15 MIN}
  \begin{columns}[T,totalwidth=\textwidth]
    \begin{column}{0.32\textwidth}
      \DMLCard{DMLBlue}{1. Contexto clínico}{Problema, evidencia publicada y verticales iniciales.\\[0.12cm]{\scriptsize 5 min}}
    \end{column}
    \begin{column}{0.32\textwidth}
      \DMLCard{DMLGreen}{2. Plataforma y prototipo}{Solución, arquitectura, módulos actuales y despliegue clínico.\\[0.12cm]{\scriptsize 6 min}}
    \end{column}
    \begin{column}{0.32\textwidth}
      \DMLCard{DMLGold}{3. Ruta estratégica}{Encaje BSC, validación y próximos 90 días.\\[0.12cm]{\scriptsize 4 min}}
    \end{column}
  \end{columns}
  \vspace{0.35cm}
  \DMLBox{\textbf{Digital Machine Listening:} convierte sonidos clínicos —voz, pulmón y corazón— en datos trazables, comparables y preparados para investigación, seguimiento e IA.}
\end{frame}

% 3
\begin{frame}{Problema clínico}
  \begin{columns}[T,totalwidth=\textwidth]
    \begin{column}{0.50\textwidth}
      \DMLEyebrow{AUDIO CON VALOR, DATOS SIN ESTRUCTURA}
      \begin{itemize}
        \item Voz, tos, respiración y auscultación contienen información longitudinal.
        \item En la práctica, la captura suele depender del dispositivo, ruido, distancia y operador.
        \item Muchos audios quedan como archivos aislados, sin metadatos ni trazabilidad.
        \item Sin calidad y contexto, es difícil construir evidencia o entrenar modelos confiables.
      \end{itemize}
    \end{column}
    \begin{column}{0.46\textwidth}
      \DMLImage{05_health_audio_triade.png}{Voz, pulmón y corazón conectados por ondas de sonido.}
    \end{column}
  \end{columns}
\end{frame}

% 4
\begin{frame}{Evidencia que abre la oportunidad}
  \DMLEyebrow{SEÑALES CLÍNICAS AUDIBLES}
  \begin{columns}[T,totalwidth=\textwidth]
    \begin{column}{0.32\textwidth}
      \DMLCard{DMLGold}{Corazón}{En un estudio en atención primaria, la auscultación tuvo sensibilidad de 44\% para enfermedad valvular significativa y 32\% para enfermedad leve.\textsuperscript{[1]}}
    \end{column}
    \begin{column}{0.32\textwidth}
      \DMLCard{DMLGreen}{Pulmón}{Un metaanálisis de análisis computarizado de sonidos pulmonares reportó sensibilidad de 80\% y especificidad de 85\% para detectar sibilancias o crepitantes.\textsuperscript{[2]}}
    \end{column}
    \begin{column}{0.32\textwidth}
      \DMLCard{DMLBlue}{Voz}{Revisiones en Parkinson muestran cambios acústicos medibles, incluyendo jitter y shimmer, y describen la voz como un posible fenotipo digital.\textsuperscript{[3,4]}}
    \end{column}
  \end{columns}
  \vspace{0.30cm}
  \DMLBox{La oportunidad no es sustituir al clínico: es capturar mejor, medir de forma reproducible y crear datos longitudinales que hoy se pierden.}
\end{frame}

% 5
\begin{frame}{La oportunidad clínica}
  \DMLEyebrow{TRES VERTICALES INICIALES}
  \begin{columns}[T,totalwidth=\textwidth]
    \begin{column}{0.32\textwidth}
      \DMLCard{DMLBlue}{Voz}{Disfonía, fatiga vocal, seguimiento de patologías, habla espontánea, lectura guiada y datos en español.\textsuperscript{[3,4]}}
    \end{column}
    \begin{column}{0.32\textwidth}
      \DMLCard{DMLGreen}{Pulmón}{Tos, respiración, sibilancias, crepitantes, calidad de auscultación y seguimiento respiratorio.\textsuperscript{[2,5]}}
    \end{column}
    \begin{column}{0.32\textwidth}
      \DMLCard{DMLGold}{Corazón}{Soplos, ritmo, calidad de señal, auscultación digital y apoyo a investigación clínica.\textsuperscript{[1,6]}}
    \end{column}
  \end{columns}
  \vspace{0.32cm}
  \DMLBox{Primer posicionamiento recomendado: investigación, seguimiento y soporte no diagnóstico, evitando afirmaciones médicas prematuras.}
\end{frame}

% 5
\begin{frame}{Solución propuesta}
  \DMLEyebrow{PLATAFORMA CLÍNICA MODULAR}
  \begin{columns}[T,totalwidth=\textwidth]
    \begin{column}{0.32\textwidth}
      \DMLCard{DMLGreen}{Captura guiada}{Protocolos, consentimiento, dispositivo, contexto clínico, ruido y control de calidad.}
    \end{column}
    \begin{column}{0.32\textwidth}
      \DMLCard{DMLBlue}{Análisis reproducible}{Productos sobre una base común: métricas, artefactos, tablas, gráficas e histórico.}
    \end{column}
    \begin{column}{0.32\textwidth}
      \DMLCard{DMLGold}{Datos para IA}{Datasets comparables, metadatos, trazabilidad, gobernanza y preparación para modelos futuros.}
    \end{column}
  \end{columns}
  \vspace{0.28cm}
  \DMLBox{La idea no es vender una prueba aislada: es construir la capa de infraestructura para productos clínicos basados en señales audibles.}
\end{frame}

% 6
\begin{frame}{Lo construido hasta ahora}
  \DMLEyebrow{PROTOTIPO FUNCIONAL}
  \begin{columns}[T,totalwidth=\textwidth]
    \begin{column}{0.31\textwidth}
      \DMLMetric{DMLGreen}{Portal}{Login, sesiones, usuarios, organizaciones y roles por producto.}
    \end{column}
    \begin{column}{0.31\textwidth}
      \DMLMetric{DMLBlue}{2 módulos}{Audioprint y Audiometer ya funcionan como productos.}
    \end{column}
    \begin{column}{0.31\textwidth}
      \DMLMetric{DMLGold}{Cobro base}{Coins por producto, saldo por usuario y ledger auditable.}
    \end{column}
  \end{columns}
  \vspace{0.28cm}
  \begin{columns}[T,totalwidth=\textwidth]
    \begin{column}{0.48\textwidth}
      \DMLCard{DMLBlue}{Audioprint}{Entrada general para analizar audio, extraer features, conservar artefactos y construir histórico.}
    \end{column}
    \begin{column}{0.48\textwidth}
      \DMLCard{DMLGold}{Audiometer}{Screening auditivo orientativo con tonos, reporte visual, histórico y consumo controlado.}
    \end{column}
  \end{columns}
\end{frame}

% 7
\begin{frame}{Arquitectura del flujo clínico}
  \begin{columns}[T,totalwidth=\textwidth]
    \begin{column}{0.52\textwidth}
      \DMLEyebrow{DE CAPTURA A EVIDENCIA}
      \begin{enumerate}
        \item Paciente, contexto, consentimiento y organización.
        \item Captura o subida de audio con metadatos.
        \item Control de calidad y repetición si la señal no sirve.
        \item Análisis por producto clínico.
        \item Reporte, histórico y comparación longitudinal.
        \item Dataset gobernado para investigación/modelos.
      \end{enumerate}
    \end{column}
    \begin{column}{0.44\textwidth}
      \DMLImage{04_platform_architecture.png}{Diagrama de plataforma clínica modular.}
    \end{column}
  \end{columns}
\end{frame}

% 8
\begin{frame}{Siguiente producto: asistente de voz clínico}
  \begin{columns}[T,totalwidth=\textwidth]
    \begin{column}{0.52\textwidth}
      \DMLEyebrow{CAPTURA CLÍNICA EN TIEMPO REAL}
      \begin{itemize}
        \item Guía protocolos de voz, tos, respiración o fonación.
        \item Valida ruido, intensidad, duración y distancia al micrófono.
        \item Extrae marcadores acústicos: tono, energía, pausas, prosodia, formantes y estabilidad.
        \item Permite seguimiento longitudinal y datasets comparables en español.
      \end{itemize}
    \end{column}
    \begin{column}{0.44\textwidth}
      \DMLImage{06_voice_assistant_clinical.png}{Asistente de voz clínico capturando audio en tiempo real.}
    \end{column}
  \end{columns}
\end{frame}

% 9
\begin{frame}{Modelo de despliegue clínico}
  \DMLEyebrow{PILOTOS CONTROLADOS Y ESCALABLES}
  \begin{columns}[T,totalwidth=\textwidth]
    \begin{column}{0.32\textwidth}
      \DMLCard{DMLGreen}{Organizaciones}{Clínicas, servicios, laboratorios o grupos de investigación con usuarios y admins propios.}
    \end{column}
    \begin{column}{0.32\textwidth}
      \DMLCard{DMLBlue}{Coins por uso}{Cada análisis consume créditos por producto: facilita pilotos, control de coste y futura compra.}
    \end{column}
    \begin{column}{0.32\textwidth}
      \DMLCard{DMLGold}{Gobernanza}{Separación por contexto, histórico, trazabilidad y borrado de datos asociados al usuario.}
    \end{column}
  \end{columns}
  \vspace{0.30cm}
  \DMLBox{La plataforma ya está pensada para pilotos multi-organización y para evolucionar hacia planes, compra de créditos y servicios de análisis.}
\end{frame}

% 10
\begin{frame}{Una plataforma, otras aplicaciones}
  \DMLEyebrow{FUERA DEL FOCO INICIAL}
  \begin{columns}[T,totalwidth=\textwidth]
    \begin{column}{0.45\textwidth}
      \DMLImage{03_audible_signals_map.png}{Mapa de señales audibles aplicables más allá de salud.}
    \end{column}
    \begin{column}{0.51\textwidth}
      \begin{itemize}
        \item Monitorización de máquinas y mantenimiento predictivo.
        \item Audio ambiental y eventos acústicos.
        \item Seguridad operacional y detección de anomalías sonoras.
        \item Calidad de procesos donde el sonido es una huella medible.
      \end{itemize}
      \vspace{0.20cm}
      \DMLSmallBox{\textbf{Alcance:} estas aplicaciones demuestran amplitud, pero la presentación y el primer negocio deben concentrarse en salud.}
    \end{column}
  \end{columns}
\end{frame}

% 11
\begin{frame}{Por que puede encajar con BSC}
  \DMLEyebrow{IA, SALUD DIGITAL, DATOS Y TRANSFERENCIA}
  \begin{columns}[T,totalwidth=\textwidth]
    \begin{column}{0.32\textwidth}
      \DMLCard{DMLBlue}{IA + HPC}{Pipelines reproducibles, entrenamiento, inferencia, escalado y análisis multimodal.}
    \end{column}
    \begin{column}{0.32\textwidth}
      \DMLCard{DMLGreen}{Rigor clínico}{Protocolos, colaboraciones, privacidad, sesgos, validación y evaluación responsable.}
    \end{column}
    \begin{column}{0.32\textwidth}
      \DMLCard{DMLGold}{Transferencia}{Transformar prototipo en activo tecnológico: pilotos, IP, partners y modelo de explotación.}
    \end{column}
  \end{columns}
  \vspace{0.32cm}
  \DMLBox{La pregunta no es solo técnica: qué ruta transforma el prototipo en evidencia clínica, producto y oportunidad real.}
\end{frame}

% 12
\begin{frame}{Próxima decisión estratégica}
  \begin{columns}[T,totalwidth=\textwidth]
    \begin{column}{0.52\textwidth}
      \DMLEyebrow{DECISIONES CRÍTICAS}
      \begin{enumerate}
        \item ¿Qué vertical clínica debería priorizarse primero?
        \item ¿Qué piloto mínimo daría evidencia en 90 días?
        \item ¿Qué colaboradores clínicos o internos deberían verlo?
        \item ¿Qué ruta parece más viable: BSC, transferencia o spin-off?
        \item ¿Qué no debería construirse antes de validar?
      \end{enumerate}
    \end{column}
    \begin{column}{0.44\textwidth}
      \DMLEyebrow{PROPUESTA 90 DÍAS}
      \begin{itemize}
        \item Elegir una vertical: voz, pulmón o corazón.
        \item Definir protocolo y consentimiento.
        \item Ejecutar piloto pequeño con organización controlada.
        \item Producir dataset inicial con gobernanza clara.
        \item Decidir ruta institucional o emprendedora.
      \end{itemize}
    \end{column}
  \end{columns}
\end{frame}

\begin{frame}{Referencias seleccionadas}
  \tiny
  \begin{enumerate}
    \item Gardezi SKM et al. \textit{Cardiac auscultation poorly predicts the presence of valvular heart disease in asymptomatic primary care patients}. Heart. 2018. doi:10.1136/heartjnl-2018-313082.
    \item Gurung A et al. \textit{Computerized lung sound analysis as diagnostic aid for the detection of abnormal lung sounds: a systematic review and meta-analysis}. Respiratory Medicine. 2011. doi:10.1016/j.rmed.2011.05.007.
    \item Chiaramonte R, Bonfiglio M. \textit{Acoustic analysis of voice in Parkinson's disease: a systematic review and meta-analysis}. Revista de Neurología. 2020. doi:10.33588/rn.7011.2019414.
    \item Tracy JM et al. \textit{Investigating voice as a biomarker: Deep phenotyping methods for early detection of Parkinson's disease}. Journal of Biomedical Informatics. 2020. doi:10.1016/j.jbi.2019.103362.
    \item Landry V et al. \textit{Audio-based digital biomarkers in diagnosing and managing respiratory diseases: a systematic review and bibliometric analysis}. European Respiratory Review. 2025. doi:10.1183/16000617.0246-2024.
    \item Chorba JS et al. \textit{Deep learning algorithm for automated cardiac murmur detection via a digital stethoscope platform}. Journal of the American Heart Association. 2021. doi:10.1161/JAHA.120.019905.
  \end{enumerate}
  \vspace{0.25cm}
  {\scriptsize Estas referencias orientan la oportunidad clínica; d-ml se plantea inicialmente como plataforma de investigación, seguimiento y soporte no diagnóstico.}
\end{frame}

\end{document}
